\documentclass[10pt,a4paper]{amsart}
\usepackage[margin=19mm]{geometry}
\usepackage{amsmath,amssymb,mathtools}
\usepackage[scr=rsfs,cal=euler]{mathalfa}
\usepackage{xfrac}
\usepackage[unicode,hidelinks,bookmarks=false]{hyperref}
\newcommand{\ie}{i.e.\ }
\newcommand{\cf}{cf.\ }
\newcommand{\resp}{respectively\ }
\newcommand{\Subsection}{\subsection}
\providecommand{\nobreakdash}{\nobreak\hbox{-}}
\setlength{\emergencystretch}{2em}
\multlinegap0pt
\usepackage{amssymb}
\usepackage{enumerate}
\newtheorem{lemma}{Lemma}
\title{Commutation Properties of Semi-groups for Contact Type Hamilton-Jacobi Equation}
\author{Guyu Jin}
\date{Source: arXiv:2603.19650v1 (CC BY 4.0)}
\begin{document}
\maketitle

\noindent\emph{Source and licence.} Commutation Properties of Semi-groups for Contact Type Hamilton-Jacobi Equation. Version arXiv:2603.19650v1 (CC BY 4.0), distributed by arXiv under the Creative Commons Attribution 4.0 International licence, http://creativecommons.org/licenses/by/4.0/, as recorded in the arXiv licence metadata for this exact version and shown on the abstract page https://arxiv.org/abs/2603.19650v1. Full licence notice: \texttt{licenses/arxiv-2603.19650-CC-BY-4.0.txt}.\par\smallskip

\noindent\textit{Editorial scope.} Counted unit: the article's lemma lemma (source0.tex lines 247--249). The article's complete original proof of it is delivered with it (source0.tex lines 250--273, one \texttt{proof} environment of 24 lines). No mathematics is written, repaired or re-derived: the statement and the proof are the article's own lines, in the article's own notation, and no retained line is substituted or re-flowed.  No further source-local context is needed: every cross-reference of every retained line resolves inside the two delivered spans.  Nothing was omitted from any retained span, so the package carries no omission record.  No retained line contains a citation, so the wrapper carries no bibliography and no bibliography entry.  No retained line contains a figure environment, an image-inclusion command or drawing code, so no image is delivered and the package declares no asset dependency.  Editorial formatting only, in addition: an AMS \texttt{amsart} preamble that restates the article's own declarations and macros used by the retained lines, namely the environments \texttt{lemma} on separate counters, together with the standard packages those lines need.  The retained environments print in this document as Lemma 1 in order of appearance, whereas the article prints the same items with the numbering of its own full document.  The complete native source of the article as distributed in the arXiv e-print for this exact version is bundled unchanged as \texttt{sources/196792/source0.tex}.
\begin{lemma}
    $u(x,t)$ is a variational solution of (HJ).
\end{lemma}

\begin{proof}
    We need to check that $u$ satisfies (i) and (ii) in Definition 1.3.
    
    (i)Let $\gamma:[t_1,t_2] \to \mathbb R^n$ be a continuous piecewise curve $C^1$ curve. Let $\bar \gamma:[0,t_1]\to \mathbb R^n$be the minimizer of $u(\gamma(t_1),t_1)$. Consider a curve $\xi :[0,t_2]\to \mathbb R^n$ defined by
    \begin{equation}
        \xi(t) = 
        \begin{cases} 
        \tilde{\gamma}(t), & t \in [0, t_1], \\
        \gamma(t), & t \in (t_1, t_2].
        \end{cases}
    \end{equation}
    It follows that
    \begin{align*}
        &u(\gamma(t_2),t_2)-u(\gamma(t_1),t_1)\\
        &=\underset{\gamma_2(t_2)=\gamma(t_2)}{\inf}\{u_0(\gamma_2(0))+\int_0^{t_2}H^*(\gamma_2(\tau),\overset{\cdot}{\gamma_2 }(\tau),u(\gamma_2(\tau),\tau))d\tau\}\\
        &-\underset{\gamma_1(t_1)=\gamma(t_1)}{\inf}\{u_0(\gamma_1(0))+\int_0^{t_1}H^*(\gamma_1(\tau),\overset{\cdot}{\gamma_1 }(\tau),u(\gamma_1(\tau),\tau))d\tau\}\\
        &\leq u_0(\xi(0))+\int_0^{t_2}H^*(\xi(\tau),\overset{\cdot}{\xi }(\tau),u(\xi(\tau),\tau))d\tau-u_0(\bar \gamma(0))-\int_0^{t_2}H^*(\bar \gamma(\tau),\overset{\cdot}{\bar \gamma }(\tau),u(\bar \gamma(\tau),\tau))d\tau,
    \end{align*}
    which implies that 
    \begin{equation}
        u(\gamma(t_2),t_2)-u(\gamma(t_1),t_1)\leq \int _{t_1}^{t_2}H^*(\gamma(\tau),\overset{\cdot}{\gamma}(\tau),u(\gamma(\tau),\tau))d\tau;
    \end{equation}
    (ii) follows from the Markov property for $u$, which is easy to show.
\end{proof}

\end{document}
